\documentclass[11pt,a4paper]{amsart}

\usepackage[T1]{fontenc}
\usepackage[utf8]{inputenc}
\usepackage{amsmath,amssymb,amsthm}
\usepackage{tikz}
\usepackage[hidelinks]{hyperref}
\usepackage{booktabs}

\theoremstyle{plain}
\newtheorem{theorem}{Theorem}[section]
\newtheorem{lemma}[theorem]{Lemma}
\newtheorem{proposition}[theorem]{Proposition}
\newtheorem{corollary}[theorem]{Corollary}
\theoremstyle{definition}
\newtheorem{definition}[theorem]{Definition}
\newtheorem{convention}[theorem]{Convention}
\newtheorem{remark}[theorem]{Remark}


\begin{document}

\title[Three conjectures on alternating plane graphs]
      {Three conjectures on alternating plane graphs}

\author{Hanyu Yang}
\email{hanyu.yang.92@gmail.com}
\urladdr{\href{https://orcid.org/0009-0005-0419-4070}{ORCID 0009-0005-0419-4070}}

\date{\today}

\subjclass[2020]{05C10, 05C30, 68R10}
\keywords{alternating plane graph, plane graph, rotation system,
computer-assisted proof}

\begin{abstract}
In an alternating plane graph, adjacent vertices have different degrees and
adjacent faces have different sizes. Four problems on these graphs were left
open by Alth\"ofer, Haugland, Scherer, Schneider and Van Cleemput. We settle two
and advance a third.

No alternating plane graph has exactly two degrees or exactly two face sizes
(Conjecture~10.1). One per-edge inequality, applied twice, proves this, and the
result survives when bridges are allowed.

Conjecture~10.2, that $(3,4,5)$-alternating plane graphs exist for every
$n \ge 20$, is settled. We give verified certificates at the $26$ orders that
were open, and a periodic capping lemma reaching order $48$ and every
$n \ge 50$. Together these close every $n \ge 46$ independently of the original
paper.

For Conjecture~10.3 we give $3$-connected witnesses at $54$ orders. Its
infinite tail is open, and so is the fourth problem.
\end{abstract}

\maketitle

\section{Introduction}

An alternating plane graph is a plane graph in which neighbours differ: adjacent
vertices have different degrees, and adjacent faces have different sizes.
Alth\"ofer, Haugland, Scherer, Schneider and Van Cleemput~\cite{AHSSV2015}
introduced these graphs and closed their paper with four open problems. We
settle the first two and advance the third.

Conjecture~10.1 says there are no $2,Y$-alternating and no $X,2$-alternating
plane graphs, where an $X,Y$-alternating plane graph has exactly $X$ distinct
degrees and exactly $Y$ distinct face sizes. Alth\"ofer et al.\ showed that a
$2,Y$-alternating plane graph has degrees $3$ and $4$, or $3$ and $5$, with at
least $25$ or $56$ vertices respectively. For weak
alternating plane graphs, which allow degree $2$, they bounded the number of
faces by summing $1/r+1/s$ over edges between an $r$-face and an $s$-face.

That face count never uses degree $2$, and it is most of the proof. Add the same
count over the two ends of each edge. By Euler's formula the four reciprocals,
summed over all edges, give $e+2$, so some edge carries more than $1$. Alternation alone allows
$\frac13+\frac14$ on each side, which is $\frac76$ and too much. But two
values on one side make the graph bipartite there, so the values on the other
side are even. Distinct even values are at least $4$ and $6$, and
$\frac14+\frac16=\frac5{12}$. Now $\frac7{12}+\frac5{12}=1$ at every edge, a
contradiction (Theorem~\ref{thm:101}). In \S\ref{sec:c2} we allow bridges. A
bridge sits on a face of size at least $8$, and that pays for the few edges that
can exceed $1$.

Conjecture~10.2 asks for a $(3,4,5)$-alternating plane graph on every
$n \ge 20$. Twenty-six orders were open, and we give certificates for all of
them. They are capped pieces of one periodic strip, and the capping lemma
(Theorem~\ref{thm:pumping}) shows that inserting further periods keeps the graph
alternating, because every face lies inside five consecutive periods.

Conjecture~10.3 asks for $3$-connected alternating plane graphs. We give
witnesses at $54$ orders (Theorem~\ref{thm:103}). The spliced graphs above
order $56$ are not known to be $3$-connected, and $(3,4,5)$-alternating plane
graphs need not be (Proposition~\ref{prop:not3conn}). The asymptotic degree
distribution stays open (\S\ref{sec:density}).

The proofs in \S\S\ref{sec:identity}--\ref{sec:c2} use no computer. The
certificates run to $110$ vertices. We give them as rotation systems and
recompute every property from these. The capping lemma rests on finite checks
of a single certificate, listed in Remark~\ref{rem:hypotheses}.

\section{Preliminaries}\label{sec:prelim}

\begin{definition}[\cite{AHSSV2015}, Definition~2.1]\label{def:apg}
An \emph{alternating plane graph} is a simple connected plane graph in
which
\begin{enumerate}
  \item[(a)] every vertex has degree at least $3$;
  \item[(b)] every face, the unbounded one included, has size at least $3$;
  \item[(c)] adjacent vertices have different degrees;
  \item[(d)] adjacent faces have different sizes.
\end{enumerate}
\end{definition}

\begin{definition}[\cite{AHSSV2015}, Definitions~3.1 and~3.4]
A $(3,4,5)$-alternating plane graph is one in which every vertex degree and
every face size lies in $\{3,4,5\}$. An \emph{$X,Y$-alternating plane graph} is
one with exactly $X$ distinct vertex degrees and exactly $Y$ distinct face
sizes.
\end{definition}

\begin{convention}[face size]\label{conv:c1}
The \emph{size} of a face is the length of its boundary walk counted with
multiplicity~\cite[\S3.3]{MoharThomassen}, that is, the number of edge-side
incidences. A bridge lies on one face from both sides and contributes $2$ to its
size.
\end{convention}

In a bipartite plane graph every closed walk is even, so every face has even
size under this convention. Counting distinct boundary vertices would lose that.
This is also the count behind $\sum_s s\,f_s = 2e$ in (3.1) and (9.2)
of~\cite{AHSSV2015}.

\begin{convention}[bridges]\label{conv:c2}
Say that a plane graph satisfies \emph{(C2)} if no edge has the same face on
both sides, equivalently if it has no bridge. In (d), the faces on the two sides
of an edge must have different sizes, so a bridge breaks (d) and alternating
plane graphs satisfy (C2)~\cite{AHSSV2015}. Under the \emph{weak reading}, (d)
constrains only pairs of distinct faces and bridges are permitted.
\end{convention}

Section~\ref{sec:101} uses (C2), and \S\ref{sec:c2} proves the same result under
the weak reading.

\begin{theorem}[\cite{AHSSV2015}, Theorem~3.2]\label{thm:32}
In a $(3,4,5)$-alternating plane graph, $v_3 = f_3$, $v_4 = f_4$ and
$v_5 = f_5$, where $v_i$ and $f_i$ count vertices of degree $i$ and faces of
size $i$. Consequently $v_5 = v_3 - 4$, $E = 2V - 2$ and $F = V$.
\end{theorem}

\subsection{Rotation systems}

A plane map is encoded by its \emph{rotation system}~\cite{MoharThomassen}: for
each vertex, its neighbours in clockwise order. It determines the map up to
homeomorphism. Write $\alpha$ for the involution exchanging the two darts of an
edge and $\sigma$ for the rotation. The faces are the orbits of
$\varphi = \sigma^{-1}\alpha$, and the size of a face is the length of its
orbit. An edge is a bridge if and only if its two darts lie in the same orbit.

\section{The per-edge identity}\label{sec:identity}

\begin{lemma}\label{lem:identity}
Let $G$ be a finite plane graph with no isolated vertex, with $v$ vertices,
$e$ edges and $f$ faces. For an edge
$a$ with ends $u,w$ and face occurrences $F^-(a)$, $F^+(a)$ on its two sides,
put
\[
  c(a) \;=\; \frac{1}{\deg u} + \frac{1}{\deg w}
           + \frac{1}{|F^-(a)|} + \frac{1}{|F^+(a)|}.
\]
Then $\sum_{a} c(a) = v + f$. If $G$ is connected, Euler's formula gives
\[
  \sum_{a} c(a) \;=\; e + 2 ,
  \qquad\text{equivalently}\qquad
  \sum_{a} \bigl(c(a) - 1\bigr) \;=\; 2 .
\]
\end{lemma}

\begin{proof}
Each vertex hands $1/\deg$ to each edge end at it, and each face hands $1/|F|$
to each edge side on its boundary walk (Convention~\ref{conv:c1}). Each pays out
exactly $1$, so the edges receive $v+f$.
\end{proof}

Neither count needs facial walks to be simple. A bridge has $F^-(a) = F^+(a)$
and receives $2/|F|$ from that face, so the lemma holds with bridges too, and
\S\ref{sec:c2} uses it that way.

We write $x(a) = c(a) - 1$ for the \emph{excess} of an edge. So
Lemma~\ref{lem:identity} reads $\sum_a x(a) = 2$, and in a connected plane graph
some edge has positive excess.

\section{Conjecture 10.1}\label{sec:101}

\begin{theorem}\label{thm:101}
There is no $2,Y$-alternating plane graph and no $X,2$-alternating plane graph.
\end{theorem}

The two halves exchange vertices and faces, but no dual graph is used, so
neither needs the $3$-edge-connectivity of~\cite[Lemma~2.2]{AHSSV2015}. We prove
both.

\subsection{No $2,Y$-alternating plane graph}

Suppose $G$ is one, with degrees $d_1 < d_2$. By (c) every edge joins a
$d_1$-vertex to a $d_2$-vertex. So $G$ is bipartite, and since $d_1 \ge 3$ the
vertex side of $c(a)$ is at most
\[
  \frac{1}{d_1} + \frac{1}{d_2} \;\le\; \frac13 + \frac14 \;=\; \frac{7}{12}.
\]
Every closed walk in $G$ is even, so every face has even size, at least $4$. By
(C2) the two sides of an edge are distinct faces, and by (d) their sizes differ.
So one is at least $4$ and the other at least $6$, and the face side is at most
$\tfrac14 + \tfrac16 = \tfrac{5}{12}$. Hence $c(a) \le 1$ for every edge,
against Lemma~\ref{lem:identity}. \qed

\begin{remark}
This is (9.5) of \cite{AHSSV2015} with $1/2$ replaced by $1/d_1 \le 1/3$. In
\S9.1 Alth\"ofer et al.\ run the same chain on the weak class with degrees $2$
and $k$. With $\tfrac12$ kept, it excludes exactly $k \ge 12$, which is their
(9.10). It does not exclude $k \le 10$, where they exhibit weak
$2,k$-alternating plane graphs, and it leaves $k = 11$ open, as they do without
their Lemma~9.4.
\end{remark}

\subsection{No $X,2$-alternating plane graph}\label{sec:halftwo}

Suppose $G$ is one, with face sizes $s_1 < s_2$. By (C2) and (d) every edge
separates an $s_1$-face from an $s_2$-face, so the face side is at most
$\tfrac13 + \tfrac14 = \tfrac{7}{12}$.

Every degree is even. Fix a vertex $u$ of degree $d$, let its edges in rotation
order be $a_1,\dots,a_d$, and let $F_i$ be the face in the corner between $a_i$
and $a_{i+1}$, indices modulo $d$. The two sides of $a_i$ are $F_{i-1}$ and
$F_i$, distinct faces by (C2), so $|F_{i-1}| \ne |F_i|$. With only two sizes
the cyclic word $|F_1|,\dots,|F_d|$ alternates, so $d$ is even.

Degrees are at least $3$ and even, hence at least $4$, and adjacent degrees
differ. So the vertex side is at most $\tfrac14+\tfrac16 = \tfrac{5}{12}$, and
again $c(a) \le 1$ for every edge. \qed

\section{Removing the dependence on (C2)}\label{sec:c2}

Throughout this section we take the weak reading of Convention~\ref{conv:c2}.
A bridge breaks both halves of the argument. Its two sides are one face, and its
ends need not have even degree. What rescues it is that a bridge sits on a face
of size at least $8$. So every bridge has excess at most $-\tfrac16$, while a
positive edge has excess at most $\tfrac1{24}$ and has a bridge at its
degree-$3$ end. Each bridge takes at most two positive edges, and the bridges
win.

\begin{lemma}[parity]\label{lem:parity}
Let $G$ be a plane graph whose faces admit a colouring by two colours that is
proper on every pair of distinct adjacent faces. Then for every vertex $v$,
\[
  \deg v \;-\; \bigl|\{\text{bridges at } v\}\bigr| \quad\text{is even.}
\]
\end{lemma}

\begin{proof}
Take the corner faces $F_1,\dots,F_d$ at $v$ as in \S\ref{sec:halftwo}. The
occurrences $F_{i-1}$ and $F_i$ are the same face exactly when $a_i$ is a
bridge. So going once round $v$, the colour changes across each non-bridge edge
and stays the same across each bridge. We return to the first colour, so the
number of changes is even.
\end{proof}

With no bridges, every degree is even. In an $X,2$-alternating plane graph the
face size gives the colouring.

\begin{corollary}\label{cor:evendeg}
In an $X,2$-alternating plane graph under the weak reading, a vertex lying on no
bridge has even degree, hence degree at least $4$.
\end{corollary}

\begin{theorem}\label{thm:nobridge}
Under the weak reading there is no $X,2$-alternating plane graph. In particular
no such graph has a bridge.
\end{theorem}

\begin{proof}
Let $G$ be one, with face sizes $s_1 < s_2$, and let $B$ be its number of
bridges. Suppose first $B \ge 1$.

\smallskip
\emph{Step 0: $s_2 \ge 8$.} Let $a = uw$ be a bridge, and let $A \ni u$ and
$C \ni w$ be the components of $G - a$. Then $\deg_A u \ge 2$ and
$\deg_C w \ge 2$. Let $W_A$ be the boundary walk of $A$ through the corner left
by deleting $a$. It cannot have length $1$ (a loop) or $2$ (parallel edges, or
$A = K_2$, but $\deg_A u \ge 2$). So $|W_A| \ge 3$, and likewise
$|W_C| \ge 3$. The face $F$ of $G$ at $a$ runs along $W_A$, across $a$, along
$W_C$ and back across $a$, so
\[
  |F| \;=\; |W_A| + |W_C| + 2 \;\ge\; 3 + 3 + 2 \;=\; 8 .
\]
As $|F| \in \{s_1,s_2\}$, we get $s_2 \ge 8$.

\smallskip
\emph{Step 1: only an edge at a vertex of degree $3$ can have positive excess.}
We have $s_1 \ge 3$, $s_2 \ge 8$, and the two ends of an edge have distinct
degrees. Three cases:
\begin{itemize}
  \item $a$ a bridge: both sides carry $F$, so
        $c(a) \le \tfrac13+\tfrac14+\tfrac28 = \tfrac56$ and
        $x(a) \le -\tfrac16$;
  \item $a$ not a bridge, both ends of degree at least $4$:
        $c(a) \le \tfrac14+\tfrac15+\tfrac13+\tfrac18 = \tfrac{109}{120}$ and
        $x(a) \le -\tfrac{11}{120}$;
  \item $a$ not a bridge with an end of degree $3$:
        $c(a) \le \tfrac13+\tfrac14+\tfrac13+\tfrac18 = \tfrac{25}{24}$ and
        $x(a) \le \tfrac1{24}$.
\end{itemize}
Only the third class can be positive.

\smallskip
\emph{Step 2: each bridge carries at most two positive edges.} This is the hard
step. A positive edge is a $3$--$4$ edge: its face side is at most
$\tfrac13+\tfrac18 = \tfrac{11}{24}$, so its vertex side exceeds
$\tfrac{13}{24}$, and $\tfrac13+\tfrac15 = \tfrac{12.8}{24}$ already falls short.

Let $a$ be a positive edge and $u$ its end of degree $3$. By
Lemma~\ref{lem:parity} the number of bridges at $u$ is odd. It is not $3$,
since $a$ is not a bridge. So $u$ lies on exactly one bridge $b(u)$, and we
charge $a$ to $b(u)$.

Now fix a bridge $b$. Every positive edge charged to $b$ has its degree-$3$ end
on $b$. By (c) at most one end of $b$ has degree $3$, and at that end only two
edges other than $b$ remain. Writing $P$ for the number of positive edges,
$P \le 2B$.

\smallskip
\emph{Step 3.} Every other edge has $x(a) \le 0$. So by Steps~1 and~2 and
Lemma~\ref{lem:identity},
\[
  2 \;=\; \sum_a x(a) \;\le\; \frac{P}{24} - \frac{B}{6}
     \;\le\; \frac{2B}{24} - \frac{B}{6} \;=\; -\frac{B}{12} \;<\; 0 ,
\]
a contradiction.

\smallskip
Finally suppose $B = 0$. By Corollary~\ref{cor:evendeg} every degree is even and
at least $4$, so the vertex side of every edge is at most
$\tfrac14+\tfrac16 = \tfrac5{12}$ and the face side at most
$\tfrac13+\tfrac14 = \tfrac7{12}$. Hence $x(a) \le 0$ for every edge, against
$\sum_a x(a) = 2$.
\end{proof}

\begin{remark}
Step~2 has room to spare. Without the degree condition at the ends of a bridge
we get only $P \le 4B$. The sum in Step~3 is then at most $0$, which still
contradicts $2$. At $P \le 5B$ the argument fails. Step~0 has room too: with
$s_2 \ge 7$ a bridge and its two positive edges total at most $-\tfrac1{84}$,
but with $s_2 \ge 6$ the bound is $+\tfrac1{12}$.
\end{remark}

\begin{proposition}\label{prop:twoY-weak}
Under the weak reading there is no $2,Y$-alternating plane graph either.
\end{proposition}

\begin{proof}
Step~0 above uses only minimum degree $3$ and simplicity, so a bridge's face has
size at least $8$. As in \S\ref{sec:101} the graph is bipartite, so every closed
walk is even and every face has even size at least $4$. A bridge has
$c(a) \le \tfrac13 + \tfrac14 + \tfrac28 = \tfrac56$. A non-bridge edge has two
distinct faces of distinct even sizes on its sides, so
$c(a) \le \tfrac13+\tfrac14+\tfrac14+\tfrac16 = 1$. So $x(a) \le 0$ for every
edge, against $\sum_a x(a) = 2$.
\end{proof}

\begin{remark}
Weak-reading alternating plane graphs with bridges exist. Join three copies of a
$(3,4,5)$-alternating plane graph by two bridges drawn in the exterior, at a
degree-$5$ vertex of each. The joined degrees become $7,6,6$, their old
neighbours keep degree $3$ or $4$, and the merged exterior face has size at
least $13$, so it differs from every interior face. This graph has neither
exactly two degrees nor exactly two face sizes.
\end{remark}

\section{The periodic capping lemma}\label{sec:pumping}

Conjecture~10.2 asks for a $(3,4,5)$-alternating plane graph on $n$ vertices for
every $n \ge 20$. Twenty-six orders were open. We supply certificates for all of
them (\S\ref{sec:artifact}), and prove that a single certificate generates an
infinite family.

\subsection{The cover}

The certificates are capped unrollings of a periodic quotient of the torus. The
relevant structure is a bi-infinite \emph{strip}, built from identical
\emph{copies}, each contributing three vertices and six edges, together with two
\emph{caps} that close it off at the ends. Write $\mathrm{TARGET}_n$ for a
certificate of order $n$. Its alignment against the cover exhibits a maximal run
of copies whose rotations are translates of one another; we call this run the
\emph{deep block} $D$.

\begin{theorem}[periodic capping lemma]\label{thm:pumping}
Fix an order $n$ whose alignment against the cover has a deep block $D$ of at
least five copies, and let $\mathrm{cut}$ be a copy at distance at least two
from each end of $D$. Then for every integer $d \ge -(|D| - 1)$ the rotation
system $\mathrm{splice}(n,d)$, obtained from $\mathrm{TARGET}_n$ by inserting
$d$ further translates of the copy $\mathrm{cut}$ (or deleting $-d$ of them when
$d < 0$), is a $(3,4,5)$-alternating plane graph on $n + 3d$ vertices.
\end{theorem}

\begin{proof}
Write $S = \mathrm{splice}(n,d)$. Every vertex of $S$ gets its rotation in one
of four ways. Cap vertices keep theirs, with strip neighbours above the cut
shifted by $d$. Strip copies at or below the cut keep theirs. A strip copy
$c$ above $\mathrm{cut}+d$ takes the rotation of copy $c-d$. The $d$ fresh
copies take the rotation of $\mathrm{cut}$, translated.

The argument is local. A facial walk is traced by
$\varphi = \sigma^{-1}\alpha$, which reads only the edge partner and the
rotation predecessor at each vertex it visits.

We bound face spans in $\mathrm{TARGET}_n$, since the faces of $S$ are what is
in question. In $\mathrm{TARGET}_n$ every face has at most five darts, and the
cover's largest edge offset is two. A closed sequence of at most five steps,
each of size at most two, has span at most four: reaching level $5$ needs three
ascents, leaving two steps to descend five. So \emph{every face of
$\mathrm{TARGET}_n$ lies inside five consecutive copies}. Both inputs are finite
checks on the one certificate (Remark~\ref{rem:hypotheses}). The bound four is
not sharp; the largest span attained, over every spliceable order at
$d = 1,2,3,5$, is two.

The fresh copies lie in $(\mathrm{cut}, \mathrm{cut}+d]$ and are translates of
$\mathrm{cut}$. Since $\mathrm{cut}$ is at distance at least two from each end
of $D$, the copies $\mathrm{cut}-2,\dots,\mathrm{cut}+2$ lie in $D$ and are
mutual translates. Hence every copy in $[\mathrm{cut}-2, \mathrm{cut}+d+2]$ is a
translate of $\mathrm{cut}$, and every five-copy window meeting the fresh block
is a translate of $[\mathrm{cut}-2,\mathrm{cut}+2]$. Away from that range $S$
agrees with $\mathrm{TARGET}_n$.

Now trace $\varphi$ in $S$ from any dart. Its counterpart in
$\mathrm{TARGET}_n$ closes within five steps inside a five-copy window, on
which $S$ is a translate of $\mathrm{TARGET}_n$. So the walk in $S$ closes at
the same length, and \emph{every face of $S$ is a translated face of
$\mathrm{TARGET}_n$}.

Each condition of Definition~\ref{def:apg} is local, so $S$ inherits it: face
sizes and degrees lie in $\{3,4,5\}$, adjacent vertices have unequal degrees and
adjacent faces unequal sizes, and facial walks are simple. A loop or parallel
edge would sit inside a two-copy window, where $S$ and $\mathrm{TARGET}_n$
agree. $S$ is connected, because each copy meets the next and every cap vertex
reaches the strip as before.

Finally $S$ is spherical. A copy contributes three vertices, six edges and,
since every face is a translated face, three faces. As $3-6+3=0$, $V-E+F$ is
$2$, as in $\mathrm{TARGET}_n$.
\end{proof}

\begin{remark}[what is machine-checked]\label{rem:hypotheses}
Three facts in the proof are checked by computation on the single certificate
$\mathrm{TARGET}_n$. Every face has at most five darts. No strip edge joins
copies more than two apart. The rotations in the deep block are translates of
one another. None is a statement about $S$.
\end{remark}

\subsection{What the family reaches}

The floor $-(|D|-1)$ is where deletion would consume a copy that is not deep,
that is, where the two cap collars meet; each collar is two copies wide.
Splicing changes the order by $3d$, so a certificate reaches one residue class
modulo $3$. The floors are:

\begin{center}
\begin{tabular}{ccl}
\toprule
$n \bmod 3$ & floor & family \\
\midrule
$0$ & $48$ & $48, 51, 54, 57, \dots$ \\
$1$ & $52$ & $52, 55, 58, 61, \dots$ \\
$2$ & $50$ & $50, 53, 56, 59, \dots$ \\
\bottomrule
\end{tabular}
\end{center}

\begin{corollary}\label{cor:102}
Conjecture~10.2 holds: there is a $(3,4,5)$-alternating plane graph on $n$
vertices for every $n \ge 20$.
\end{corollary}

\begin{proof}
Orders below $46$ are the published witnesses of \cite{AHSSV2015}, decoded from
the corpus~\cite{APGcorpus}, together with the Section-8 closures of
\S\ref{sec:103}. Orders $37$ and $38$ are covered by the heuristic search of
\cite{AHSSV2015}; the disk surgery of \S\ref{sec:103} gives checked witnesses
there as well. Orders $46$ through $56$, $67$ through $74$, $88$ through $92$,
and $109$, $110$ carry the certificates of this paper. Every remaining order is
reached by Theorem~\ref{thm:pumping} from a certificate in its residue class
modulo $3$, with the floors above.
\end{proof}

\begin{remark}[the deletion direction]\label{rem:deletion}
The floors in the table are reached by deletion. For $d < 0$ the locality
argument reads as for insertion, with the four cases collapsing, except at the
floor itself: there a three-copy window meets both cap collars, and not every
window is a translated window. The floor orders therefore also carry their own
verified certificates.

The reach is set by the hypothesis $|D| \ge 5$. Of the twenty spliceable
orders, five, namely $48$, $51$, $53$, $54$ and $56$, have deep blocks of only
$1, 2, 2, 3$ and $3$ copies, so they are not eligible seeds. The eligible ones
begin at $67$, $68$ and $69$, one per residue class, so insertion reaches
nothing below $70$, and all ten orders $57$--$66$ lie outside
Theorem~\ref{thm:pumping}. None of them has a stored certificate. They are
covered by the computation in \texttt{test\_pumping\_splice.py}, which builds
successive splices from $90$, $109$ and $110$ and runs each through the decision
procedures. A separate proof of the deletion direction, with its own cut rule,
is open.
\end{remark}

This gives an independent construction of (\cite{AHSSV2015}, Theorem~8.1), which covers
$n \ge 111$, and of twenty-three of the twenty-six open orders, from one
certificate per residue class. Orders $46$, $47$ and $49$ lie outside the family
and rest on their certificates alone. So the certificates and
Theorem~\ref{thm:pumping} settle \emph{every order $n \ge 46$} without
Theorem~8.1 of~\cite{AHSSV2015}, subject to Remark~\ref{rem:deletion} at the
five orders noted there. Orders $20$ through $45$ are taken from that paper.

\begin{remark}
The splice reproduces certificates it did not start from.
$\mathrm{TARGET}_{110}$ shortened by twenty periods is $\mathrm{TARGET}_{50}$ up
to isomorphism of rotation systems, and so is $\mathrm{TARGET}_{92}$ shortened
by fourteen; $\mathrm{TARGET}_{90}$ shortened by fourteen and
$\mathrm{TARGET}_{72}$ shortened by eight both give $\mathrm{TARGET}_{48}$. So
certificates found independently decompose into the same pair of caps and the
same period.
\end{remark}

\section{Conjecture 10.3}\label{sec:103}

\begin{theorem}\label{thm:103}
There is a $3$-connected alternating plane graph on $n$ vertices for
$$n \in \{17\} \cup [19,56] \cup [67,74] \cup [88,92] \cup \{109,110\},$$
$54$ orders in all, each certified by an explicit witness. (Order $17$ is
below the range of Conjecture~10.3; it is included because it is the smallest
order at which any alternating plane graph exists, and its witness is
$3$-connected.)
\end{theorem}

Conjecture~10.3 asks for every $n \ge 19$. The remaining orders $57$--$66$,
$75$--$87$, $93$--$108$ and $n \ge 111$ are reached only by the spliced family,
and we cannot prove that family $3$-connected (Remark~\ref{rem:103gap}).

The witnesses come from six sources, and no order depends on more than one of
them being correct.

\begin{center}
\begin{tabular}{ll}
\toprule
orders & source \\
\midrule
$19$ & all five alternating plane graphs on $19$ vertices \\
$17, 20, 26$--$36, 42$ & published witnesses, decoded and re-verified \\
$21, 22, 23, 25$ & published search seeds \\
$21$--$24$, $39$--$45$ & Section-8 closures, rebuilt here \\
$37, 38$ & disk surgery \\
$46$--$56$, $67$--$74$, $88$--$92$, $109$, $110$ & the certificates of \S\ref{sec:pumping} \\
\bottomrule
\end{tabular}
\end{center}

\emph{Order $19$.} Conjecture~10.2 cannot help here: \cite{AHSSV2015} reports
no $(3,4,5)$-alternating plane graph on $18$ or $19$ vertices, so a witness must
be a general alternating plane graph. There are exactly five on
$19$ vertices~\cite{HoG}, and all five are $3$-connected.

\emph{Orders $37$ and $38$.} The heuristic search of \cite{AHSSV2015} reaches
every order from $20$ to $42$, but its Section-8 construction, of order
$18a+19b+20c+21d+3$, misses these two, because $34$ and $35$ are not sums of
$18,19,20,21$. So the witnesses built here are the first checkable from a stored
rotation system. We obtain them by disk surgery: excise the star of a vertex, an
edge or a face from a stored $(3,4,5)$-alternating plane graph, let the boundary
degrees float, and refill the disk.

\subsection{Connectivity is not automatic}

If every $(3,4,5)$-alternating plane graph were $3$-connected, then
Corollary~\ref{cor:102} would give Theorem~\ref{thm:103} for all $n \ge 20$ in
one step.

\begin{proposition}\label{prop:not3conn}
There is a $(3,4,5)$-alternating plane graph on $46$ vertices with a separating
pair. In particular $(3,4,5)$-alternating plane graphs need not be
$3$-connected.
\end{proposition}

The graph is given as a rotation system in the artifact. Both verifiers accept
it, its separating pair is unique, and removing that pair leaves two components
of $22$ vertices each. Theorem~\ref{thm:103} does not depend on the proposition,
since every witness is checked individually. But it means the connectivity of
the spliced family needs a proof.

\begin{remark}[the infinite tail]\label{rem:103gap}
The natural induction shows that a freshly spliced copy is adjacent to the
copies on both sides, and concludes that connectivity is preserved. What is
needed is that $S(n,d) - \{u,v\}$ is connected, with the pair already removed,
and that does not follow. The induction also needs the splice point to be placed
clear of a pair occupying up to four consecutive copies, inside a deep block of
at least five, and we have not checked this.

The witnesses at $17$, $19$--$56$, $67$--$74$, $88$--$92$ and $109$--$110$ do
not depend on the induction. Orders $57$--$66$, $75$--$87$, $93$--$108$ and
every $n \ge 111$ do. \cite{AHSSV2015} states in its concluding remarks that the
$(3,4,5)$-alternating plane graphs of its Section~8 are $3$-connected, which
would cover exactly those gaps, but gives no proof. At those orders
Conjecture~10.3 is open.
\end{remark}

\section{The asymptotic degree distribution}\label{sec:density}

The fourth problem of~\cite[\S10]{AHSSV2015} asks about the shape of a typical
graph. Writing $r = v_3$, Theorem~\ref{thm:32} confines $v_4$ to
$[r-5, \tfrac32 r - 5]$, so the ratio $v_4/v_3$ lies asymptotically in
$[1, 1.5]$. Alth\"ofer et al.\ ask whether there is a density function on that
interval giving the asymptotic fractions of $(3,4,5)$-alternating plane graphs
for large $n$, and if so what it looks like. We cannot answer it.

All twenty-six certificates satisfy $v_5 = v_3 - 4$ and lie at the bottom of
the interval. Their $v_4/v_3$ ranges over $[1.0000, 1.1250]$ and drifts downwards
with $n$, reaching exactly $1$ at orders $50$, $56$, $74$, $92$ and $110$. This
reflects the construction. The certificates are capped unrollings of one
periodic strip, and Theorem~\ref{thm:pumping} inserts copies of a single period,
so every member of the family inherits that period's degree profile.

Our methods do not reach the question. Everything proved here is an existence
statement, while a density needs a measure over all graphs of order $n$. The
class is also rigid: no two-edge switch of any of the twenty-six certificates
leaves a valid $(3,4,5)$-alternating plane graph, so local moves cannot walk the
space.

We see three routes. An exhaustive census at small orders, with
\texttt{plantri}~\cite{plantri} or a similar generator restricted to the degree
and face-size profile, would give a first histogram of $v_4/v_3$ and suggest a
conjecture. A structure theorem giving \emph{every} $(3,4,5)$-alternating plane
graph a cap--strip--cap decomposition with a bounded interface alphabet would
make the count a regular language and the density a transfer-matrix
computation; we have only a construction. Failing both, a
bijection with an already-counted family of plane maps would serve, and the
identities of Theorem~\ref{thm:32} are the natural place to look for one.

\section{The artifact}\label{sec:artifact}

Every graph in this paper is given as a rotation system in the artifact, in a
documented JSON format with a dependency-free reference reader. No file records
that a graph is alternating. Degrees, faces, face sizes, bridges, connectivity
and both alternation conditions are recomputed on every run by four separately
written decision procedures. Three are for the $(3,4,5)$ class. The fourth is
for the general class of Definition~\ref{def:apg}, which Conjecture~10.3 needs
at order~$19$. Certificates can also be exported to \texttt{planar\_code} for
checking by \texttt{plantri}~\cite{plantri} or House of Graphs~\cite{HoG}.

The artifact runs $1308$ checks. The decision procedures use only the Python
standard library, and the test suite also needs \texttt{pytest}. Each
certificate is accepted by every verifier of its class, and the order-$19$
witnesses are rejected by the three $(3,4,5)$ procedures. The computations the
proofs rest on are the finite hypotheses of Remark~\ref{rem:hypotheses}.

\subsection*{Data availability}

The artifact, with every certificate, the four decision procedures, a drawing
tool (a straight-line embedding after Tutte~\cite{Tutte1963}), the export tools
and the test suite, is archived at~\cite{artifact} under the MIT licence. The
concept DOI \href{https://doi.org/10.5281/zenodo.22269200}{10.5281/zenodo.22269200}
resolves to the current version; this paper describes version~$1.0.4$,
\href{https://doi.org/10.5281/zenodo.22666407}{10.5281/zenodo.22666407}.
Reproduce with \texttt{make deps \&\& make verify}. The published
corpus~\cite{APGcorpus} carries no licence statement, so its graphs are
re-encoded in the artifact's format, with the SHA-256 of each original kept, and
the MIT licence does not cover them.

\section*{Disclosure of AI use}

Large language models were used throughout this work. They helped search for
arguments and reviewed the proofs adversarially. The decision procedures and
the search programs that found the certificates were written with AI
assistance. The prose was drafted by AI from the author's notes, and the author
edited it. The idea of dropping (C2) (\S\ref{sec:c2}) came from an AI
reviewer. The systems were Claude Opus~5 and Claude Fable~5.1 (Anthropic),
ChatGPT at its Pro tier and Codex \texttt{gpt-5.6-terra} and
\texttt{gpt-5.6-sol} (OpenAI), and DeepSeek \texttt{v4-pro} and
\texttt{v4-flash}, used from February to September 2026. The author checked the
mathematics and is responsible for it. No AI system is an author.

\begin{thebibliography}{9}

\bibitem{AHSSV2015}
I.~Alth\"ofer, J.~K. Haugland, K.~Scherer, F.~Schneider and N.~Van Cleemput,
\emph{Alternating plane graphs},
Ars Math. Contemp. \textbf{8} (2015), 337--363.
\href{https://doi.org/10.26493/1855-3974.584.09a}{doi:10.26493/1855-3974.584.09a}.

\bibitem{artifact}
H.~Yang, \emph{Alternating plane graphs: settling Conjectures 10.1 and 10.2,
with witnesses for 10.3} --- certificates, verifiers and test suite, version 1.0.5, Zenodo, 2026.
\href{https://doi.org/10.5281/zenodo.22269200}{doi:10.5281/zenodo.22269200}
(concept).

\bibitem{MoharThomassen}
B.~Mohar and C.~Thomassen, \emph{Graphs on Surfaces},
Johns Hopkins University Press, Baltimore, 2001.

\bibitem{Tutte1963}
W.~T. Tutte, \emph{How to draw a graph},
Proc. London Math. Soc. (3) \textbf{13} (1963), 743--767.
\href{https://doi.org/10.1112/plms/s3-13.1.743}{doi:10.1112/plms/s3-13.1.743}.

\bibitem{plantri}
G.~Brinkmann and B.~D. McKay,
\emph{Fast generation of planar graphs} (expanded edition),
MATCH Commun. Math. Comput. Chem. \textbf{58} (2007), 323--357.

\bibitem{HoG}
K.~Coolsaet, S.~D'hondt and J.~Goedgebeur,
\emph{House of Graphs 2.0: A database of interesting graphs and more},
Discrete Appl. Math. \textbf{325} (2023), 97--107.
\href{https://doi.org/10.1016/j.dam.2022.10.013}{doi:10.1016/j.dam.2022.10.013}.

\bibitem{APGcorpus}
I.~Alth\"ofer, \emph{Alternating plane graphs} --- status page and
machine-readable corpus of published witnesses,
\url{https://althofer.de/alternating-plane-graphs.html} and
\url{https://althofer.de/apg/apgs/}, accessed 1 September 2026.

\end{thebibliography}

\end{document}
